% !TeX program = xelatex
% Самостоятельный документ в кодировке UTF-8.
\documentclass[12pt,a4paper]{article}
\usepackage{fontspec}
\usepackage{unicode-math}
\setmainfont{DejaVu Serif}
\setmonofont{DejaVu Sans Mono}[Scale=MatchLowercase]
\usepackage{babel}
\babelprovide[import,main]{russian}
\usepackage[margin=25mm]{geometry}
\usepackage[hidelinks]{hyperref}
\setlength{\emergencystretch}{2em}
\newcommand{\parag}[1]{\S\,#1}
\newcommand{\code}[1]{\texttt{#1}}
\newcommand{\task}[1]{\par\medskip\noindent\textbf{#1}\quad}
\title{План СНИРС\\«Экспериментальное сравнение сканеров»}
\author{Дисциплина «Грамматика и автоматы»}
\date{}
\begin{document}
\maketitle
\raggedright

\noindent\textbf{Основание работы.} В \parag{2.6.1} «Интерпретаторы таблиц» и \parag{2.6.2} «Представление автомата в программном коде» рассматриваются два способа реализации переходов конечного автомата.

\begin{quote}
Как зависит время выделения лексем от способа программного представления функции переходов конечного автомата? Как эта зависимость проявляется на входных словах разной длины и разного лексического состава?
\end{quote}

В первом опыте сравнить три сканера, решающих одну лексическую задачу. Во втором опыте у собственного сканера менять только способ вычисления функции переходов $\delta$. Первый опыт показывает различия между реализациями в целом. Второй позволяет исследовать различие, описанное в \parag{2.6.1} и \parag{2.6.2}.

\section{Какие готовые сканеры взять и какую задачу им поставить}

Взять два генератора сканеров, flex и re2c. По одному описанию лексики подготовить правила для каждого генератора, получить две программы на C и сравнивать работу этих программ.
\begin{itemize}
\item \textbf{flex.} Сгенерировать сканер в режиме \code{-Cf} с полными таблицами. Это практический пример хранения переходов в таблицах, связанный с \parag{2.6.1}. См.\ \href{https://westes.github.io/flex/manual/Options-for-Scanner-Speed-and-Size.html}{руководство flex, раздел о представлении таблиц}.
\item \textbf{re2c.} По тем же лексическим правилам сгенерировать сканер с переходами, представленными в программном коде. Это практический пример подхода из \parag{2.6.2}. См.\ \href{https://re2c.org/}{описание генератора re2c}.
\end{itemize}

\task{Общая задача.} Взять фрагмент псевдокода из \parag{0.2.4} «Псевдокод». Согласно \parag{3.2.1} «Структура описания языка программирования», определить алфавит, затем лексику --- множество лексических единиц. Классы лексем описать регулярными выражениями в смысле \parag{1.2.1} «Регулярные операции» и \parag{1.2.2} «Алгебра регулярных выражений».
\begin{itemize}
\item \textbf{Идентификаторы и числа.} Воспользоваться правилами из \parag{0.2.4}. Идентификатор начинается латинской буквой, за которой могут следовать буквы и цифры. Число состоит из одной или нескольких цифр.
\item \textbf{Ключевые слова.} Выбрать слова из псевдокода, например \code{if}, \code{then}, \code{else}, \code{while}, \code{do}, \code{return}. Окончательный перечень закрепить после отбора учебных фрагментов.
\item \textbf{Знаки операций и разделители.} Перечислить их явно и зафиксировать текстовую запись математических обозначений, например \code{:=} для присваивания.
\item \textbf{Пропуски.} Пробелы, табуляции, переводы строк и комментарии распознавать и пропускать. Комментарий начинается с \code{//} и продолжается до конца строки. Предусмотреть комментарий в конце файла.
\item \textbf{Выход.} Для каждой выданной лексемы записывать её вид, начальную позицию и позицию сразу за лексемой. При ошибке выдавать начальную позицию первого фрагмента, из которого нельзя выделить лексему.
\end{itemize}

В текущей позиции выбирать самое длинное подходящее слово. При равной длине применять заранее установленный приоритет. Например, правило ключевого слова \code{if} имеет приоритет перед правилом идентификатора. Тогда \code{if} --- ключевое слово, \code{if2} --- один идентификатор, а \code{if:=15} --- три лексемы. Такое соглашение согласуется с \href{https://westes.github.io/flex/manual/Matching.html}{правилами выбора совпадения во flex}.

\task{Результат пункта 1.} Единое описание лексики и два построенных по нему сканера. Для любого входа определено, какой результат они должны выдать.

\section{Какой сканер написать самостоятельно}

Написать на C сканер на основе детерминированного конечного автомата. Исходная модель --- автомат-распознаватель из \parag{2.2.1}. Задание таблицей рассматривается в \parag{2.1.2}, графом переходов состояний --- в \parag{2.1.3}. Образец табличного распознавателя приведён в \parag{2.6.1}.

\begin{enumerate}
\item \textbf{Построить граф переходов.} Указать входной алфавит $A$, множество состояний $Q$, начальное состояние $q_0$, заключительные состояния $F$ и функцию переходов $\delta$. Каждому заключительному состоянию сопоставить вид лексемы либо пропуск. Разобрать на графе слова \code{if} и \code{if2}. Для небольшого фрагмента можно сразу построить детерминированный граф. При построении через недетерминированный источник использовать детерминизацию из \parag{1.2.4}.
\item \textbf{Записать тот же граф таблицей.} Строка соответствует состоянию, столбец --- входному символу либо группе символов с одинаковыми переходами. Обозначить случай, когда очередной символ уже не может продолжить лексему.
\item \textbf{Дополнить распознаватель выделением границ.} Распознаватель из \parag{2.6.1} получает слово целиком. Сканеру его конец заранее неизвестен. При чтении запоминать последнюю позицию, в которой достигнуто заключительное состояние, и соответствующий вид лексемы. Когда продолжение невозможно либо вход закончился, выдать самое длинное найденное слово или исполнить пропуск. Следующую лексему начинать с первого символа, не вошедшего в предыдущую. Если ни одно заключительное состояние не достигнуто, зафиксировать ошибку.
\item \textbf{Подготовить два варианта вычисления $\delta$.} В первом следующий шаг определяется обращением к таблице $\delta[q,a]$, во втором --- ветвлениями по состоянию и входному символу, как в \parag{2.6.2}. Оба варианта получить из одного описания переходов. Чтение входа, начальное состояние, выбор границ и обработка результата общие.
\end{enumerate}

\task{Результат пункта 2.} Один собственный сканер с двумя вариантами вычисления $\delta$. Вместе с программами, построенными с помощью flex и re2c, получаются четыре измеряемых варианта. Собственный табличный вариант участвует в обоих опытах.

\section{Как провести вычислительный эксперимент}

\task{Опыт А. Сравнение сканеров.} Сравнить flex \code{-Cf}, re2c и собственный табличный сканер. Зафиксировать лексические правила, тексты, требуемый результат, язык C, компилятор, параметры сборки и машину. Установить, как различается время работы трёх реализаций на каждом семействе входов.

\task{Опыт Б. Сравнение представлений автомата.} Сравнить два варианта собственного сканера по \parag{2.6.1} и \parag{2.6.2}. Граф, функция $\delta$, входные тексты и код выделения лексем общие. Измерить, как меняется время при замене табличного вычисления $\delta$ ветвлениями в программном коде.

\begin{enumerate}
\item \textbf{Проверить корректность.} Сравнить виды и границы всех лексем на каждом входе. Для коротких примеров выписать ответы вручную. Для искусственных входов сохранять ожидаемый результат при построении текста. На ошибочных входах сверить позицию ошибки. До измерения времени устранить расхождения.
\item \textbf{Определить измеряемый участок.} Заранее загрузить текст в память и подготовить буфер. Измерять проход от начала текста до конца. Генерацию программы, компиляцию, чтение с диска, подготовку буфера и печать вынести за пределы измерения. Во время прохода одинаково учитывать результат, например вычислять общую для всех реализаций контрольную сумму по видам и границам лексем.
\item \textbf{Повторить измерения.} Выполнить, например, 15 повторов для каждой пары «вариант --- вход», чередуя порядок запуска вариантов. Для коротких текстов измерять серию полных проходов и делить время на их число.
\item \textbf{Записать показатели.} Для каждого входа указать длину $N$ в символах и число выданных лексем $K$. Для времени $T$ привести медиану и разброс повторов. Основной показатель --- время на входной символ $T/N$, дополнительный --- число выданных лексем в секунду $K/T$.
\item \textbf{Сделать вывод по каждому опыту.} Опыт А показывает, какая реализация быстрее на данном семействе при данных условиях. Вывод о влиянии способа вычисления $\delta$ основывать на опыте Б. Если разница сопоставима с разбросом повторов, указать, что в выбранных условиях устойчивого преимущества обнаружить не удалось.
\end{enumerate}

\section{Как составить представительный набор нагрузочных тестов}

Обосновать представительность относительно выбранного фрагмента псевдокода. Охватить все заданные виды лексем, частое и редкое завершение лексем, общие начала слов, пропуск комментариев и рост длины всего входа. Использовать учебные тексты и управляемые искусственные серии.

\begin{enumerate}
\item \textbf{Алгоритмы из учебника.} Найти в исходном файле учебника вхождения \verb|\begin{algorithmic}|. Взять, например, интерпретатор из \parag{2.6.1}, схему ветвлений из \parag{2.6.2} и алгоритм построения функции переходов из \parag{2.6.6}. Привести математические обозначения к выбранной текстовой записи и сохранить список замен. Для каждого образца указать объём, частоты видов лексем, их длины и долю пропусков. Это позволит описать различия учебных текстов с точки зрения сканера.
\item \textbf{Короткие и длинные лексемы.} При одинаковом общем объёме построить тексты из идентификаторов длины 2, 8, 32 и 128, разделённых одним пробелом. Повторить для чисел. Сопоставить время на символ с числом завершений лексем и возвратов в начальное состояние.
\item \textbf{Общие начала ключевых слов и идентификаторов.} Сравнить парные входы с \code{ifx} и \code{abx}, затем с \code{while2} и \code{abcde2}. Сохранить длину, число идентификаторов, их позиции и разделители. В следующих входах менять долю идентификаторов, начинающихся с ключевого слова. Виды и границы лексем совпадают. По графу показать, какие состояния посещает собственный автомат.
\item \textbf{Комментарии и пробельные символы.} При фиксированном объёме менять долю текста внутри комментариев. Отдельно сравнить много коротких и несколько длинных комментариев при одинаковом суммарном объёме комментариев. Сканер читает эти символы и выполняет переходы, хотя в выдаче лексем им ничего не соответствует. Сопоставить $T/N$ и $K/T$.
\item \textbf{Рост длины входного слова.} Для фиксированной смеси видов и длин лексем построить входы, например, на 64~КиБ, 1~МиБ и 16~МиБ. Зафиксировать способ генерации и начальные значения генератора. Исследовать изменение $T$ и $T/N$ при увеличении $N$ и сохранении лексического состава.
\end{enumerate}

\task{Короткие проверочные примеры.} Подготовить пустой вход, \code{if}, \code{if2}, \code{if:=15}, последнюю лексему без пробела, комментарий без завершающего перевода строки и недопустимый символ после правильной лексемы. Для каждого случая заранее выписать ожидаемые виды и границы лексем либо позицию ошибки.

\task{Дополнительный опыт при наличии времени.} В конце \parag{2.6.1} отмечено, что табличное представление эффективно, если массивы переходов «не слишком велики». Зафиксировать входной текст и постепенно добавлять в описание лексики отсутствующие в нём ключевые слова. Для каждой версии записывать число состояний, размер таблицы в байтах и время. Отдельным прогоном посчитать число различных состояний, фактически посещённых на данном входе. При увеличении таблицы сканер может обращаться почти к тем же строкам.

\task{Порядок выполнения.} Зафиксировать лексику по \parag{0.2.4} и \parag{3.2.1}; построить и описать автомат по \parag{2.1} и \parag{2.2}; получить реализации по \parag{2.6.1} и \parag{2.6.2}, а также сканеры flex и re2c; проверить результаты на коротких примерах; подготовить пять семейств нагрузок; провести опыты А и Б. Для каждого вывода назвать семейство входных слов, измеренную величину и свойство реализации автомата, к которому относится вывод.

\section*{Распределение работы между двумя участниками}

\task{Совместное начало.} Выбрать учебные фрагменты. Согласовать алфавит, перечень и приоритет правил, обозначения видов лексем, единицу и начало отсчёта позиций, поведение на ошибке и конце входа. Утвердить общий интерфейс подачи текста и выдачи результата. Вручную разобрать короткие проверочные примеры и сохранить ответы.

\task{Участник 1. Собственный автомат и опыт Б.} Построить граф и таблицу переходов. Из одного описания получить табличный вариант и вариант с ветвлениями. Отвечать за код собственного сканера и примеры путей по графу. Подготовить учебные тексты с перечнем замен. По общим измерениям разобрать результаты опыта Б.

\task{Участник 2. Готовые сканеры и опыт А.} Подготовить правила flex и re2c и подключить обе программы к общему интерфейсу. Построить четыре искусственных семейства входов с ожидаемыми ответами. Отвечать за программу проверки и измерений, сбор результатов и их обработку. По общим измерениям разобрать результаты опыта А.

\task{Совместная сверка.} Каждый редактирует файлы своего участка. Общую спецификацию ведёт участник 1, программу измерений --- участник 2. Изменения правил согласовываются до внесения в реализации. Сначала подключить все четыре варианта и сверить результаты на коротких примерах, затем проверить весь набор входов. Итоговые замеры выполнить одной программой на одной машине. Выводы проверить вместе, сопоставив данные опыта А с контрольным сравнением в опыте Б.

\end{document}
